%%%%%%%% LAYER FINGERPRINTS -- ICML DRAFT %%%%%%%%%%%%%%%%%
% REWRITTEN 2026-09-16 from scratch (previous body deleted) against PAPER_POINTS.md,
% the curated, scope-anchored story for this specific paper (see
% notion/Alex Fingerprinting/ICLR 27 Drafting *.md for the scope doc itself). Content
% and numbers are transcribed from PAPER_POINTS.md / STATUS.md / PROJECT_PROGRESS.md;
% if this draft and those ever disagree, the markdown files are the source of truth.
% Headlined by the actually-strongest result (joint necessity, replicated) rather than
% the single-layer geometry-predicts-causal-importance framing of the previous draft,
% which is a null result, not a positive finding -- see PAPER_POINTS.md's "Honest
% stock-take" section for the full ranking this structure follows.

\documentclass{article}

\usepackage{microtype}
\usepackage{graphicx}
\usepackage{subcaption}
\usepackage{booktabs}
\usepackage{hyperref}
\newtheorem{definition}{Definition}
\usepackage{array}
\usepackage{ragged2e}
\usepackage{multirow}
\usepackage{multicol}
\usepackage{enumitem}
\usepackage{makecell}
\usepackage{verbatim}
\usepackage[most]{tcolorbox}
\usepackage{parskip}

\newcommand{\theHalgorithm}{\arabic{algorithm}}

\usepackage{icml2026}

\usepackage{amsmath}
\usepackage{amssymb}
\usepackage{mathtools}
\usepackage{amsthm}

\usepackage[capitalize,noabbrev]{cleveref}

\theoremstyle{plain}
\newtheorem{theorem}{Theorem}[section]
\newtheorem{proposition}[theorem]{Proposition}
\newtheorem{lemma}[theorem]{Lemma}
\newtheorem{corollary}[theorem]{Corollary}
\theoremstyle{definition}
\newtheorem{assumption}[theorem]{Assumption}
\theoremstyle{remark}
\newtheorem{remark}[theorem]{Remark}

\usepackage[textsize=tiny]{todonotes}

\icmltitlerunning{Layer Fingerprints: Structure, Relevance, and Joint Necessity in CoT}

\begin{document}

\twocolumn[
  \icmltitle{Individually Redundant, Jointly Necessary: Label-Free Layer Fingerprints\\
             and Relevance-Filtered Structure in Chain-of-Thought Reasoning}

  \icmlsetsymbol{equal}{*}

  \begin{icmlauthorlist}
    \icmlauthor{Firstname1 Lastname1}{yyy}
    \icmlauthor{Firstname2 Lastname2}{yyy}
  \end{icmlauthorlist}

  \icmlaffiliation{yyy}{Department of XXX, University of YYY, Location, Country}
  \icmlcorrespondingauthor{Firstname1 Lastname1}{first1.last1@xxx.edu}

  \icmlkeywords{interpretability, chain-of-thought, layer-wise relevance propagation, mechanistic interpretability}

  \vskip 0.3in
]

\printAffiliationsAndNotice{}

\begin{abstract}
We study how a token's residual-stream representation \emph{changes} across depth
during chain-of-thought (CoT) generation, not just its final state, via a compressed,
label-free \emph{fingerprint}: a random-projected summary of a token's per-layer deltas
concatenated across every transition. On our primary model (Qwen3-0.6B, six synthetic
task families) this object carries real, non-random structure: per-(token,
layer-transition) clusters are 89--99\% single-transition-specific, and independent
component analysis recovers a dedicated \texttt{<think>}/\texttt{</think>} component in
every one of ten dataset/model combinations tested, including a 13$\times$ larger
second model (Qwen3-8B), where both results replicate and strengthen. Three geometric
metrics (path tortuosity, transition-specificity, and multi-lag directional persistence)
converge on a further finding: each model carries its own task-invariant geometric
signature, distinct from other models -- a genuine model-level fingerprint, not merely a
per-layer one. We then test this paper's central methodological question directly: does
restricting each token's residual-stream update to only its most LRP-relevant attention
heads (projected through their own output weights) make delta-based clustering less
noisy than pooling every head? The answer is real but asymmetric -- it improves recovery
of one token class in most cases tested and degrades another in a task-family-specific,
not fully robust way -- and we report this honestly rather than as a clean confirmation.
Separately, we identify a concrete attention-level mechanism (a specific head set, ruling
out the textbook synthetic induction head as a negative control) that implements ordinal
successor prediction and generalizes across task categories while relocating in depth.
Finally, we show this structure is causally load-bearing: three individually
low-importance, individually redundant layers are jointly necessary -- bypassing them
collapses the CoT trace's answer while reasoning stays coherent -- and this survives
escalating controls and \emph{replicates with a stronger effect} on a second,
structurally unrelated task family (41 of 41 samples, $p=9\times10^{-13}$, versus the
original 30 of 30, $p=1.9\times10^{-9}$). Single-layer geometric and causal metrics fail
to even detect this known effect, a validity problem rather than a null result,
motivating the group-level and relevance-filtered analyses above over simpler per-layer
description.
\end{abstract}

\section{Introduction}
\label{sec:intro}

Standard interpretability practice analyzes what a token's representation \emph{is} at a
given layer. A complementary signal is available: what \emph{changes} between layers --
a token's per-layer delta -- which in principle indexes the computation performed on it,
not just its resting state. If tokens undergoing similar computations produce similar
deltas, clustering delta sequences should recover functionally meaningful groups without
ever seeing a task label. Prior work on CoT (e.g. \emph{thought anchors},
\citealp{thoughtanchors2025}) argues that some reasoning steps are more important than
others, but has not linked that claim to internal computation directly. We ask three
linked questions. First, does a compressed, label-free summary of a token's depth-wise
trajectory -- a \emph{fingerprint} -- carry real, clusterable structure at all? Second,
does filtering that structure by what the model itself judges relevant (layer-wise
relevance propagation, LRP; \citealp{bach2015lrp}) change what is recoverable? Third, is
any of this structure causally load-bearing, or merely descriptive?

We answer all three on Qwen3-0.6B \citep{qwen3} across six synthetic task families
(letter and day-of-week recitation, number-letter pairing, country-capital recall, and
indirect-object identification), with a second-model check on Qwen3-8B (13$\times$
larger, same family). Our design is intentionally narrow: we do not claim state
representations are less informative than deltas -- an earlier, broader test of that
premise returned no clean answer across our six datasets (3 favor delta, 1 a coin flip,
2 favor state, including the best-powered one) -- deltas here are a tool for building
the fingerprint object, not a claim about representational superiority.

\section{Methods}
\label{sec:methods}

\textbf{The fingerprint object.} For each token, we compute the residual-stream delta
$\delta_t^\ell = h_t^{\ell+1} - h_t^\ell$ at every layer transition. Following
independently-fit Gaussian random projections \citep{bingham2001random} (one per
transition), we concatenate the projected deltas across all transitions into a single
per-token vector. This is entirely label-free: construction never sees a category label,
only token identities and raw activations.

\textbf{Clustering and validity.} We cluster fingerprints with spectral clustering
\citep{ng2001spectral} (nearest-neighbors affinity) across a $k$-sweep, and separately
cluster raw per-(token, layer-transition) deltas without concatenation, to test
transition-locality directly. Two token classes serve as label-free validity checks
throughout: \texttt{meta} (discourse-marker-like: \texttt{user}/\texttt{wants}/
\texttt{asked}) and \texttt{noise} (\texttt{\textbackslash n}, \texttt{<|endoftext|>},
\texttt{<|im\_end|>}, and model-specific terminal tokens).

\textbf{Independent component analysis.} On final-layer activations, we apply
PCA-whitening (15 components) followed by FastICA \citep{hyvarinen2000ica} and measure,
for each token type, the z-scored distinctiveness of its mean loading on every
component relative to the population.

\textbf{LRP relevance and head filtering.} For each token at each layer, we decompose
the attention block's exact contribution to the residual stream into its 16 (Qwen3-0.6B)
or 32 (Qwen3-8B) per-head parts: head $h$'s output (\texttt{hook\_z}) projected through
its own slice of $W_O$, summed over all heads recovers the model's real update exactly.
We rank heads by LRP relevance magnitude and compare clustering built from all heads
(\emph{unfiltered}) against clustering built from only the top-$K$ most relevant heads
per token per layer (\emph{filtered}), using \texttt{meta}/\texttt{noise} best-cluster
$F_1$ as the comparison metric.

\textbf{Causal necessity.} We test layers 24--26 (of 28) via hard zero-ablation
(forcing $\text{resid\_post}_\ell = \text{resid\_pre}_\ell$) both individually and as a
group, and via a matched-magnitude intervention design: bypassing a layer versus
injecting a magnitude-matched perturbation at the same position, escalating from
isotropic noise to a harder same-subspace control (another token's own observed
activity, direction-normalized).

\section{Results}
\label{sec:results}

We report results in order of evidentiary strength, following an explicit self-audit
rather than the order in which analyses were run.

\subsection{Delta-based structure is real, not noise}
\label{sec:results:structure}

Population token paths are far straighter than an isotropic random walk at the same
depth (population median tortuosity 1.84--2.02 across all six datasets, against a
$\sqrt{27}\approx5.2$ baseline) -- establishing that there is real, non-random structure
in the deltas before any clustering is attempted. Per-(token, layer-transition) clusters
are \textbf{89--93\% single-transition-specific} across a $5\times$ range of $k$
(25--120): tokens occupy a depth-local niche rather than a stable cross-layer identity.
On Qwen3-8B this \textbf{replicates and strengthens}: 95.0--99.2\% across the identical
$k$ range on two further task families, above the entire Qwen3-0.6B band at every $k$
tested. Independent component analysis recovers a dedicated \texttt{<think>}/
\texttt{</think>} component in \textbf{every one of ten dataset/model combinations}
tested (z-magnitude 4.5--13.7), including all three Qwen3-8B datasets -- the single most
consistent finding in this paper. These tokens additionally show a distinctive,
reproducible \emph{wanders-and-returns} archetype (tortuosity 2.62--2.70 against the
1.84--2.02 population baseline, a multi-lag directional persistence that never fully
decays), recovered independently of the clustering or ICA results above.

\subsection{A model-level, not merely per-layer, fingerprint}
\label{sec:results:modelfingerprint}

The per-layer clustering result above is \emph{local} (which transition a token's delta
belongs to). A separate question is whether the model as a whole carries a stable
geometric signature across tasks. Three independent metrics converge: Qwen3-0.6B's
population tortuosity sits in a tight 1.84--2.02 band across all six task families;
Qwen3-8B's own band, 1.73--1.80, is tighter still and consistently \emph{below}
Qwen3-0.6B's across all three of its datasets; Apertus-8B (a different architecture
family) lands at 2.04, above Qwen3-0.6B's range. The same within-model-consistent,
between-model-different pattern holds for transition-specificity (89--93\% vs.\
95--99\%, Qwen3-0.6B vs.\ Qwen3-8B) and multi-lag persistence peak height
(0.113--0.153 vs.\ 0.162--0.174). Only three models loaded the same way exist, one scale
change and one architecture change at once, so we cannot yet separate whether this
signature tracks model scale or architecture -- stated as an open question, not resolved
here.

\subsection{A concrete mechanism: the successor circuit}
\label{sec:results:successor}

Via direct attribution inspection (unsupervised clustering of this signal mostly failed
the same robustness check that sank two earlier, retracted subcluster claims), we find a
four-phase attribution pattern in letter-recitation traces: attention-sink $\to$
semantic anchor $\to$ \emph{immediately-preceding letter in the enumeration}
(induction-like) $\to$ local punctuation. As a negative control, we confirm layer 16,
head 14 is a genuine synthetic induction head by the standard diagnostic (score 0.974)
-- but it is \emph{not} the real driver: on every instance checked it attributes to
punctuation, not the preceding letter. The real signal is distributed across a
different head set (13, 3, 7, 9, 8; head 13 the sharpest single specialist, 53\% of
instances). This transfers to a second ordinal category (\texttt{days\_of\_week}): at
that task's own peak layer (26), the best head reaches 54.3\% (44/81), nearly identical
concentration, and is a member of the same candidate head set -- but the mechanism
\emph{relocates in depth} (16 vs.\ 26) rather than recurring at a fixed layer, consistent
with \citet{gould2024successor}'s claim of cross-category generality for this
computation type.

\subsection{Does relevance filtering help? A real but asymmetric answer}
\label{sec:results:lrp}

Across nine dataset/model combinations, filtering to the top-relevant attention heads
improves \texttt{meta}-cluster recovery in six, leaves it flat in two, and shows a
negligible decrease in one. \texttt{noise}-cluster recovery is flat in most cases but
drops substantially (F$_1$ change $-0.25$ to $-0.42$) on three datasets
(\texttt{countries\_capitals}, \texttt{ioi}, \texttt{ioi50}). A follow-up sweep across
$K\in\{2,4,6,8\}$ on those three datasets shows this degradation is \textbf{not robust}:
\texttt{countries\_capitals} improves at every $K$ except 4; \texttt{ioi50} flips sign
between $K=4$ and $K=6$; only \texttt{ioi} shows a consistent, shrinking degradation
across the range. We report this as a genuine, only partially-understood asymmetry
rather than either a clean confirmation or refutation of the relevance-filtering
hypothesis.

\subsection{Joint causal necessity: the strongest result, and it replicates}
\label{sec:results:necessity}

Hard-ablating layers 24--26 collapses the CoT trace's final answer while reasoning stays
coherent; three independent signals agree (behavioral collapse, degradation of the
otherwise-pure \texttt{noise} cluster, and a logit-lens next-token-prediction check
dropping from 89\% clean to $\sim$11\% ablated). Individually, these layers rank
21st/25th/26th of 28 by single-layer causal importance -- near the bottom, not the top --
yet bypassing all three \emph{together} collapses the answer far more than a
matched-magnitude injection control, a result that survives escalating from isotropic
noise to a harder same-subspace permutation control ($p=1.9\times10^{-9}$, 30 of 30
samples, \texttt{countries\_capitals}). This effect \textbf{replicates on a second,
structurally unrelated, well-powered task family} (\texttt{ioi50}: short, non-templated,
procedurally varied, versus \texttt{countries\_capitals}'s single-fact recall): bypass
beats the harder control in \textbf{41 of 41} measurable samples,
$p=9.1\times10^{-13}$ -- stronger on every metric than the original. This converts our
strongest positive result from a single-dataset finding into a replicated one.

\subsection{Single-layer metrics fail to detect this effect: a validity problem}
\label{sec:results:validity}

Per-layer geometric magnitude and reorientation are weak, non-significant predictors of
single-layer causal importance ($\rho=-0.207$, $p=0.29$; $\rho=0.023$, $p=0.91$;
$n=28$ layers). Because layers 24--26 individually rank near the bottom by the same
single-layer metric despite being jointly indispensable (\Cref{sec:results:necessity}),
this null is a \textbf{validity problem, not informative evidence} about whether geometry
predicts causal importance -- the instrument fails to detect an effect independently
known to exist. This motivates our use of group-level and relevance-filtered analyses
(\Cref{sec:results:lrp,sec:results:necessity}) over single-layer description throughout.

\section{Limitations}
\label{sec:limitations}

\textbf{Single model, mostly.} All results beyond \Cref{sec:results:structure,%
sec:results:modelfingerprint} (which include a Qwen3-8B check) are Qwen3-0.6B only.
Apertus-8B does not produce \texttt{<think>} output on any task tried (a trained
behavior, not a bug, confirmed on two task families) and contributes only the population
tortuosity data point above. Bridge-loaded models (architecturally incompatible with our
primary tooling) show a substantially more extreme version of the model-fingerprint
pattern (near-random-walk tortuosity); whether this is a genuine property or a
loading-path artifact is an open, unresolved question we flag rather than adjudicate.

\textbf{The relevance-filtering result needs more than one operating point.} Our
$K$-sweep (\Cref{sec:results:lrp}) covers only the three datasets that initially showed
degradation; whether the \texttt{meta}-class improvement is similarly robust across $K$
has not yet been checked.

\textbf{The founding delta-vs-state premise is unresolved, deliberately.} An early test
of whether tracking change outperforms tracking state at recovering category structure
returned no clean answer. We do not read this as invalidating the geometric or causal
results above, which describe properties of a trajectory a state-only view cannot
express by construction, but we do not claim deltas have been shown superior to state as
a general representational choice.

\section{Conclusion}
\label{sec:conclusion}

A label-free, compressed summary of how token representations change across depth
carries real, reproducible, class-specific structure, and at least one instance of that
structure -- three individually redundant layers -- is causally necessary as a group in
a way single-layer analysis cannot detect and that replicates, strengthening, across
structurally unrelated tasks. Filtering that structure by the model's own notion of
relevance changes what is recoverable, but not uniformly, and we report that asymmetry
directly rather than resolve it prematurely.

\bibliography{references}
\bibliographystyle{icml2026}

\end{document}
